\pagestyle{DS_FP}
%\NomPrenomNote{}

\bigskip

\exods{0}

\bigskip

 Le jeu du baguenaudier est un jeu de casse-tête constitué d'une réglette comportant \texttt{n} cases, numérotées de \texttt{1} à \texttt{n}. Chaque case peut être soit vide, soit contenir un pion.


 \og{} Jouer\fg{} une case consiste à placer un pion dans la case (remplir) si elle est vide ou enlever un pion (vider) si elle est remplie.
Initialement, toutes les cases sont vides.

 Le but du jeu est de remplir toutes les cases du baguenaudier en suivant les règles suivantes:

\begin{itemize}
    \item on ne peut jouer qu'une case à la fois;
    \item chaque case ne peut contenir qu'un pion;
    \item on peut toujours jouer la case \texttt{1};
    \item si le baguenaudier n'est ni vide ni rempli, on peut aussi jouer la case qui suit la première case remplie;
    \item aucune autre case ne peut être jouée.
\end{itemize}

\begin{minipage}{0.5\linewidth}
    Exemple de situation avec un baguenaudier de $5$ cases:

\end{minipage}\hfill{}
\begin{minipage}{0.5\linewidth}
\begin{center}
%\begin{figure}[!h]
\newcommand{\Rvide}[2]{\draw(#1,#2) -- ++ (0.75,0) -- ++(0,0.75) -- ++(-0.75,0) -- cycle; \draw({#1+0.375},{#2+0.375}) circle(0.15);}
\newcommand{\Rplein}[2]{\draw(#1,#2) -- ++ (0.75,0) -- ++(0,0.75) -- ++(-0.75,0) -- cycle; \draw[fill=black] ({#1+0.375},{#2+0.375}) circle(0.15);}
\centering\begin{tikzpicture}
\Rvide{0}{0}
\Rvide{0.75}{0}
\Rplein{1.5}{0}
\Rplein{2.25}{0}
\Rvide{3}{0}
\foreach \x/\val in {0.375/1,1.125/2,1.875/3,2.625/4,3.375/5}
{\draw (\x,0) node[below=0.1cm]{$\val$};}
\end{tikzpicture}
\captionof{figure}{Situation baguenaudier de $5$ cases}
\label{baguenaudier5}
%\end{figure}
\end{center}
    \end{minipage}

    \medskip


\noindent Dans cette situation, on peut poser un pion dans la case $1$ ou enlever le pion de la case $4$, mais on ne peut pas jouer les cases $2$, $3$ et $5$. En effet:
\begin{itemize}
    \item On peut jouer la case $1$ car c'est la première case du jeu.
    \item le baguenaudier n'est ni vide ni rempli, la première case remplie est la case $3$, on peut jouer la case $4$ qui suit la première case remplie.
    \item on ne peut pas jouer la case $2$ car elle ne suit pas la première case remplie.
    \item on ne peut pas jouer la case $3$ car elle ne suit pas la première case remplie.
    \item on ne peut pas jouer la case $5$ car elle ne suit pas la première case remplie.   
\end{itemize}

\bigskip

\begin{enumerate}[resume]
\item On considère un baguenaudier à $4$ cases sur lequel les trois dernières cases sont remplies. Désigner les cases que l'on peut jouer et dessiner l'état du baguenaudier à la suite de chacun des coups possibles.

    \begin{blocvisible}[height=2]
        Les cases que l'on peut jouer sont la case $1$ et la case $2$. 
        En jouant la case $1$, on obtient un baguenaudier avec les cases $1$, $3$ et $4$ remplies. En jouant la case $2$, on obtient un baguenaudier avec les cases $2$, $3$ et $4$ remplies.
    \end{blocvisible}


\end{enumerate}

\medskip

 Pour modéliser le baguenaudier on utilise un tableau (de type \verb|list|) de booléens. Une case vide est représentée par \texttt{False} et une case remplie par \texttt{True}.



 L'exemple de situation du baguenaudier de $5$ cases de la Figure \ref{baguenaudier5} présentée plus haut sera ainsi représenté par le tableau suivant: 

    \begin{CodePythontexAlt}[TaillePolice=\large, Largeur=1\linewidth,Centre,Lignes=true,PremLigne=1]{}
[False, False, True, True, False]
\end{CodePythontexAlt}

 La fonction \texttt{initialiser} prend en paramètre un nombre entier \texttt{n} et elle renvoie une liste modélisant un baguenaudier vide de \texttt{n} cases. Exemple:

    \begin{CodePythontexAlt}[TaillePolice=\large, Largeur=1\linewidth,Centre,Lignes=true,PremLigne=1]{}
    >>> initialiser(3)
    [False, False, False]
\end{CodePythontexAlt}

\pagebreak

\begin{enumerate}[resume]
\item Ecrire le code de la fonction \texttt{initialiser(n)}
\end{enumerate}

\begin{blocvisible}[height=3]
    \begin{CodePythontexAlt}[TaillePolice=\large, Largeur=1\linewidth,Centre,Lignes=true,PremLigne=1]{}
    def initialiser(n):
        return [False] * n
    \end{CodePythontexAlt}
\end{blocvisible}

\medskip

 La fonction \texttt{victoire(tab)} renvoie un booléen indiquant si toutes les cases du baguenaudier sont remplies.

    \begin{CodePythontexAlt}[TaillePolice=\large, Largeur=1\linewidth,Centre,Lignes=true,PremLigne=1]{}
def victoire(tab):
    for etat_case in tab:
        if etat_case == ......:
            return ......
    return ......

    \end{CodePythontexAlt}


        \begin{enumerate}[resume]
\item Compléter les lignes 3, 4 et 5 du code de la fonction \texttt{victoire}.
\end{enumerate}

\begin{blocvisible}[height=4]
    \begin{CodePythontexAlt}[TaillePolice=\large, Largeur=1\linewidth,Centre,Lignes=true,PremLigne=1]{}
    def victoire(tab):
        for etat_case in tab:
            if etat_case == False:
                return False
        return True
    \end{CodePythontexAlt}
\end{blocvisible}

\medskip

 La fonction \verb|indice_premiere_case_occupee(tab)| prend en paramètre un tableau \texttt{tab} modélisant l'état du baguenaudier et renvoie l'indice de la première case remplie du baguenaudier, ou \texttt{None} si le baguenaudier est vide. Exemples:

    \begin{CodePythontexAlt}[TaillePolice=\large, Largeur=1\linewidth,Centre,Lignes=true,PremLigne=1]{}
>>> indice_premiere_case_occupee([False, False, True, True, False])
2
>>> indice_premiere_case_occupee([True, True, True, True, True])
0
>>> indice_premiere_case_occupee([False, False, False, False, False])
None
    \end{CodePythontexAlt}

    \begin{enumerate}[resume]
\item Ecrire le code de la fonction \verb|indice_premiere_case_occupee|
\end{enumerate}

\begin{blocvisible}[height=6]
    \begin{CodePythontexAlt}[TaillePolice=\large, Largeur=1\linewidth,Centre,Lignes=true,PremLigne=1]{}
    def indice_premiere_case_occupee(tab):
        for i in range(len(tab)):
            if tab[i] == True:
                return i
        return None
    \end{CodePythontexAlt}
\end{blocvisible}

\pagebreak


 La fonction \verb|coup_valide(tab, case)| prend en paramètres \texttt{tab} qui est une liste modélisant l'état d'un baguenaudier, et un entier \texttt{case} qui est l'indice de la case à jouer. La fonction renvoie \texttt{True} si le coup respecte les règles et \texttt{False} si ce n'est pas le cas. L'indice de la première case du jeu est \texttt{0}. Pour que le coup soit valide, il faut aussi s'assurer que l'indice est un entier positif inférieur à la longueur de \texttt{tab}.


 \begin{enumerate}[resume]
\item Ecrire le code de la fonction \verb|coup_valide(tab, case)|
\end{enumerate}

\begin{blocvisible}[height=6]
    \begin{CodePythontexAlt}[TaillePolice=\large, Largeur=1\linewidth,Centre,Lignes=true,PremLigne=1]{}
    def coup_valide(tab, case):
        if case < 0 or case >= len(tab):
            return False
        if case == 0:
            return True
        indice_premiere_case_occupee = indice_premiere_case_occupee(tab)
        if indice_premiere_case_occupee is None:
            return True
        return case == indice_premiere_case_occupee + 1
    \end{CodePythontexAlt}
\end{blocvisible}

\medskip

La fonction \verb|changer_case(tab, case)| prend en paramètres un tableau \texttt{tab} modélisant l'état du baguenaudier et un indice de case nommé \texttt{case}. Si le coup désigné par \texttt{case} est valide, la fonction renvoie le tableau modélisant le baguenaudier obtenu après avoir changé l'état de la case correspondante. Si le coup joué n'est pas valide, le baguenaudier passé en paramètre est renvoyé sans modification. 
Exemples:

\begin{CodePythontexAlt}[TaillePolice=\large, Largeur=1\linewidth,Centre,Lignes=true,PremLigne=1]{}
>>> changer_case([False, False, False], 0)
[True, False, False]
>>> changer_case([False, False, False], 1)
[False, True, False]
>>> changer_case([False, False, False], 2)
[False, False, True]
    
\end{CodePythontexAlt}
\begin{enumerate}[resume]
\item Ecrire le code de la fonction \verb|changer_case(tab, case)|
\end{enumerate}

\begin{blocvisible}[height=6]
    \begin{CodePythontexAlt}[TaillePolice=\large, Largeur=1\linewidth,Centre,Lignes=true,PremLigne=1]{}
    def changer_case(tab, case):
        if not coup_valide(tab, case):
            return tab
        tab_modifie = tab.copy()
        tab_modifie[case] = not tab_modifie[case]
        return tab_modifie
    \end{CodePythontexAlt}
\end{blocvisible}

\pagebreak



La résolution d'un baguenaudier de $3$ cases, en Figure \ref{exemplesresolutions3} et $4$ cases, en Figure \ref{exemplesresolutions4} sont données ci-dessous.
\begin{center}
\begin{figure}[!h]
\centering\begin{subfigure}{6cm}
\newcommand{\Rvide}[2]{\draw(#1,#2) -- ++ (0.75,0) -- ++(0,0.75) -- ++(-0.75,0) -- cycle; \draw({#1+0.375},{#2+0.375}) circle(0.15);}
\newcommand{\Rplein}[2]{\draw(#1,#2) -- ++ (0.75,0) -- ++(0,0.75) -- ++(-0.75,0) -- cycle; \draw[fill=black] ({#1+0.375},{#2+0.375}) circle(0.15);}
\centering\begin{tikzpicture}
\Rvide{0}{0} \Rvide{0.75}{0} \Rvide{1.5}{0}
\Rplein{0}{-1} \Rvide{0.75}{-1} \Rvide{1.5}{-1}
\Rplein{0}{-2} \Rplein{0.75}{-2} \Rvide{1.5}{-2}
\Rvide{0}{-3} \Rplein{0.75}{-3} \Rvide{1.5}{-3}
\Rvide{0}{-4} \Rplein{0.75}{-4} \Rplein{1.5}{-4}
\Rplein{0}{-5} \Rplein{0.75}{-5} \Rplein{1.5}{-5}
\end{tikzpicture}
\caption{Résolution du baguenaudier à $3$ cases}
\label{exemplesresolutions3}
\end{subfigure}
\hspace{1cm}
\begin{subfigure}{9cm}
\newcommand{\Rvide}[2]{\draw(#1,#2) -- ++ (0.75,0) -- ++(0,0.75) -- ++(-0.75,0) -- cycle; \draw({#1+0.375},{#2+0.375}) circle(0.15);}
\newcommand{\Rplein}[2]{\draw(#1,#2) -- ++ (0.75,0) -- ++(0,0.75) -- ++(-0.75,0) -- cycle; \draw[fill=black] ({#1+0.375},{#2+0.375}) circle(0.15);}
\centering\begin{tikzpicture}
\Rvide{0}{0} \Rvide{0.75}{0} \Rvide{1.5}{0} \Rvide{2.25}{0} \Rplein{3.5}{0} \Rvide{4.25}{0} \Rplein{5}{0} \Rvide{5.75}{0}
\Rplein{0}{-1} \Rvide{0.75}{-1} \Rvide{1.5}{-1} \Rvide{2.25}{-1} \Rvide{3.5}{-1} \Rvide{4.25}{-1} \Rplein{5}{-1} \Rvide{5.75}{-1}
\Rplein{0}{-2} \Rplein{0.75}{-2} \Rvide{1.5}{-2} \Rvide{2.25}{-2} \Rvide{3.5}{-2} \Rvide{4.25}{-2} \Rplein{5}{-2} \Rplein{5.75}{-2}
\Rvide{0}{-3} \Rplein{0.75}{-3} \Rvide{1.5}{-3} \Rvide{2.25}{-3} \Rplein{3.5}{-3} \Rvide{4.25}{-3} \Rplein{5}{-3} \Rplein{5.75}{-3}
\Rvide{0}{-4} \Rplein{0.75}{-4} \Rplein{1.5}{-4} \Rvide{2.25}{-4} \Rplein{3.5}{-4} \Rplein{4.25}{-4} \Rplein{5}{-4} \Rplein{5.75}{-4}
\Rplein{0}{-5} \Rplein{0.75}{-5} \Rplein{1.5}{-5} \Rvide{2.25}{-5}
\end{tikzpicture}
\caption{Résolution du baguenaudier à $4$ cases}
\label{exemplesresolutions4}
\end{subfigure}
\caption{Quelques résolutions du baguenaudier}
\end{figure}
\end{center}


Il est possible d'obtenir la solution du jeu du baguenaudier, sous forme d'affichage, en utilisant des fonctions récursives nommées \texttt{vider} et \texttt{remplir}. Ces deux fonctions prennent en paramètre un entier \texttt{n} correspondant au nombre de cases du jeu. La fonction \texttt{vider} vide un baguenaudier de \texttt{n} cases initialement remplies. La fonction \texttt{remplir} remplit un baguenaudier de \texttt{n} cases initialement vides. 
Le code de la fonction \texttt{vider}, incomplet, est donné ci-dessous. On notera que les cases sont numérotées à partir de \texttt{1} dans la suite et que l'appel \texttt{vider(0)} ne fait rien.
\begin{CodePythontexAlt}[TaillePolice=\large, Largeur=1\linewidth,Centre,Lignes=true,PremLigne=1]{}
    def vider(n):
        if n == 1:
            print('Vider case 1')
        elif n > 1:
            vider(n-2)
            print(......)
            remplir(n-2)
            vider(n-1)
\end{CodePythontexAlt}


\begin{enumerate}[resume]
\item Recopier et compléter la ligne 6 du code de la fonction \texttt{vider} pour qu'elle affiche le texte \verb|'Vider case'| suivi de la valeur de l'entier \texttt{n}.

\begin{blocvisible}[height=2]
    \begin{CodePythontexAlt}[TaillePolice=\large, Largeur=1\linewidth,Centre,Lignes=true,PremLigne=1]{}
            print('Vider case', n)
    \end{CodePythontexAlt}
\end{blocvisible}

\medskip

\item En supposant que l'appel \texttt{remplir(1)} affiche \verb|'Remplir case 1'| et que l'appel \texttt{remplir(0)} ne fait rien, donner les affichages, dans la console, produits par l'exécution de \texttt{vider(3)}.
\begin{blocvisible}[height=5]
    \begin{verbatim}
Vider case 1
Vider case 3
Remplir case 2
Vider case 2
Vider case 4
Remplir case 3
Vider case 3
\end{verbatim}
\end{blocvisible}

\pagebreak

\item En vous inspirant de la fonction \texttt{vider}, de la résolution donnée à la question précédente pour \texttt{n = 3} et des Figures \ref{exemplesresolutions3} et \ref{exemplesresolutions4}, écrire le code de la fonction récursive \texttt{remplir(n)}.

\begin{blocvisible}[height=6]
    \begin{CodePythontexAlt}[TaillePolice=\large, Largeur=1\linewidth,Centre,Lignes=true,PremLigne=1]{}
    def remplir(n):
        if n == 1:
            print('Remplir case 1')
        elif n > 1:
            vider(n-2)
            print('Remplir case', n)
            remplir(n-2)
            vider(n-1)
    \end{CodePythontexAlt}
\end{blocvisible}
\end{enumerate}